\section{Introduction}
\label{sec:intro}

% Motivation: Why is this problem important?
Federated reinforcement learning (FRL) enables multiple agents to collaboratively learn policies while preserving data locality and privacy. In this paradigm, agents interact with their local environments and periodically share knowledge through a central aggregation mechanism, avoiding the need to transmit raw interaction data. This approach is particularly valuable in domains where data privacy is critical or communication costs are prohibitive, such as mobile robotics, personalized recommendation systems, and healthcare applications.

However, extending federated learning principles to reinforcement learning presents unique challenges:



% Gap: What is missing in current approaches?
Existing federated learning approaches, particularly those employing FedAvg-style aggregation, face significant limitations when applied to deep reinforcement learning. First, deep neural network representations lose expressiveness during the linear averaging process, as the aggregation operates on weight space rather than the learned feature representations. This representation-agnostic aggregation can degrade the quality of the learned Q-functions. Second, standard federated averaging requires all participating agents to maintain identical model architectures, preventing heterogeneous deployment scenarios where agents may have different computational capabilities or environmental complexities.


% Our approach: What do we propose?
To address these limitations, we propose FedQHD, a federated Q-learning framework that employs hyperdimensional computing for function approximation. By representing Q-functions through high-dimensional vector spaces, our approach enables linear aggregation of learned models while preserving the expressiveness of the approximated value functions. The use of hyperdimensional regression provides a principled method for combining decentralized learning experiences through simple averaging operations.




% Contributions: What are the main contributions?
The main contributions of this work are:
\begin{itemize}
    \item \textbf{Aggregation-aware Q-learning}: We demonstrate that hyperdimensional function approximation maintains expressiveness under linear averaging, enabling effective aggregation of decentralized interaction data without degrading the learned representations.
    \item \textbf{Model heterogeneity support}: Our framework naturally accommodates agents with different model configurations, allowing federated learning across heterogeneous computational environments.
\end{itemize}


% Organization: How is the paper structured?
The rest of this paper is organized as follows. Section~\ref{sec:related_work} reviews related work. Section~\ref{sec:preliminaries} provides background on federated learning, Q-learning, and hyperdimensional computing. Section~\ref{sec:methodology} describes our proposed FedQHD approach. Section~\ref{sec:experiments} details our experimental setup. Section~\ref{sec:results} presents our experimental results. Section~\ref{sec:discussion} discusses the implications of our findings. Finally, Section~\ref{sec:conclusion} concludes the paper.
